\section*{Chapitre \ref{ch:b2:magmatic-series} --- Séries magmatiques et contextes tectoniques}

\begin{solution}{exo:b2:magmatic-series:1}
(a) $\ce{SiO2} = 49.60$, alcalins $2.22 + 0.51 = 2.73$~: entre 45 et 52 et sous 5,
\emph{basalte} (B). (b) 58.5 et $3.9 + 1.6 = 5.5$~: entre 57 et 63~; la limite supérieure y va
de (57, 5.9) à (63, 7.0), soit $5.9 + 1.1\times1.5/6 = 6.18$ pour 58.5~; 5.5 est sous elle~:
\emph{andésite} (O2). (c) 72.0 et $4.0 + 4.2 = 8.2$~: au-delà de la verticale à 69, au-dessus de 8
et au-dessus de la droite de (69, 8) à (77, 0), qui vaut 5.0 pour 72~: \emph{rhyolite} (R).
(d) 47.0 et $4.5 + 2.5 = 7.0$~: à 47, le champ U1 est compris entre la droite de (45, 5) à
(49.4, 7.3), à $5 + 2.3\times2/4.4 = 6.05$, et la droite de (45, 9.4) à (49.4, 7.3), à
$9.4 - 2.1\times2/4.4 = 8.45$. La roche est dans U1~: une \emph{téphrite} si sa norme contient
moins de \qty{10}{\percent} d'olivine, une
\emph{basanite} si elle en contient davantage.
\end{solution}

\begin{solution}{exo:b2:magmatic-series:2}
(a) Les oxydes totalisent 98.40 et le total avec les volatils 100.00~;
\ce{H2O} $1.40 < 2$ et \ce{CO2} $0.20 < 0.5$~: l'analyse est acceptable. (b) Facteur
$100/98.40 = 1.01626$~: \ce{SiO2} 52.44, \ce{TiO2} 1.22, \ce{Al2O3} 16.36,
\ce{FeO_T} 9.96, \ce{MgO} 6.71, \ce{CaO} 8.33, \ce{Na2O} 3.56, \ce{K2O} 1.42
(somme 100.00)~; alcalins $3.56 + 1.42 = 4.98$. (c) Avant~: 51.60 et 4.90, juste sous la
limite basalte--andésite basaltique à 52~: \emph{basalte}. Après~: 52.44 et 4.98~; à
52.44, la limite supérieure de O1, de (52, 5) à (57, 5.9), vaut $5 + 0.9\times0.44/5
= 5.08$~: \emph{andésite basaltique} (O1). Le recalcul a fait passer la roche de l'autre côté
d'une limite~: les noms proches d'une droite dépendent d'elle. (d) À 52.44, la droite bas-K vaut $0.3 +
0.045\times4.44 = 0.50$ et la droite haut-K $1.2 + 0.095\times4.44 = 1.62$~; 1.42
se situe entre les deux~: \emph{moyen-K}.
\end{solution}

\begin{solution}{exo:b2:magmatic-series:3}
(a) $Q + A + P = 90$~: $Q = 25/90 = \qty{27.8}{\percent}$ et $100\,P/(A + P) =
100\times35/65 = 53.8$~: champ 3b, un \emph{monzogranite}, un granite. La biotite et
l'hornblende n'entrent pas dans la somme QAPF. (b) $Q + A + P = 65$~: $Q = 5/65 =
\qty{7.7}{\percent}$ (entre 5 et 20) et $100\times55/60 = 91.7$ (au-dessus de 90)~: champ
10*, diorite quartzifère, gabbro quartzifère ou anorthosite quartzifère. Le \omterm{def:b2:magmatic-series:qapf}{diagramme QAPF} seul ne
permet pas de choisir entre eux~; le plagioclase et les minéraux mafiques (ici
\qty{35}{\percent} de clinopyroxène) le permettent.
\end{solution}

\begin{solution}{exo:b2:magmatic-series:4}
La droite bas-K vaut $0.3 + 0.045(\ce{SiO2} - 48)$, la droite haut-K $1.2 +
0.095(\ce{SiO2} - 48)$. (a) À 50~: 0.39 et 1.39~; $0.40$ est juste au-dessus de la droite basse~:
\emph{moyen-K}. (b) À 56~: 0.66 et 1.96~: \emph{moyen-K}. (c) À 62~: 0.93 et 2.53~;
$2.90 > 2.53$~: \emph{haut-K}. (d) À 66~: 1.11 et 2.91~; $1.00 < 1.11$~: \emph{bas-K}.
\end{solution}

\begin{solution}{exo:b2:magmatic-series:5}
(a) Basalte~: $A + F + M = 3.2 + 14.5 + 4.5 = 22.2$~; $A = 14.4$, $F = 65.3$, $M = 20.3$.
Andésite~: $5.05 + 6.1 + 3.8 = 14.95$~; $A = 33.8$, $F = 40.8$, $M = 25.4$. (b) La droite
ne dépasse jamais $F = 54.5$~; le basalte, avec $F = 65.3$, se trouve du côté de F
quelle que soit son $A$~: \emph{tholéiitique}, un ferrobasalte d'une tendance de Fenner. Pour
l'andésite, à $A = 33.8$, la droite passe entre (49.6, 35.0) et (52.5, 30.0), à $F =
49.6 + 2.9\times1.2/5 = 50.3$~; $40.8 < 50.3$~: \emph{calco-alcaline}. (c) La magnétite,
qui emporte le fer ferrique d'un \omterm{def:b2:magma-emplacement:magma}{magma} oxydé dès le début,
et empêche $F/M$ de monter aussi vite que dans une tholéiite.
\end{solution}

\begin{solution}{exo:b2:magmatic-series:6}
MORB~: $\mathrm{Nb_N} = 4/0.658 = 6.08$, $\mathrm{La_N} = 4.12/0.648 = 6.36$, rapport
0.956. \omterm{def:b2:elements-in-the-earth:morb-oib}{Basalte d'île océanique}~: $38/0.658 = 57.8$, $30.8/0.648 = 47.5$, rapport 1.215. Basalte d'arc~:
$3/0.658 = 4.56$, $6.1/0.648 = 9.41$, rapport 0.484. Seul le basalte d'arc présente une anomalie en
niobium, un creux en Nb sur un diagramme en araignée~: le niobium y est appauvri par rapport à son
voisin La~: le \omterm{def:b2:magma-emplacement:magma}{magma} d'arc a reçu des éléments de la plaque,
mais peu de niobium. Le MORB est presque
plat~; le \omterm{def:b2:elements-in-the-earth:morb-oib}{basalte d'île océanique} est plus riche en Nb qu'en La.
\end{solution}

\begin{solution}{exo:b2:magmatic-series:7}
(a) $\mathrm{Ti} = 1.6\times0.5993\times10^4 = \qty{9589}{ppm}$~; Ti/100 = 95.9, Zr
110, $3\mathrm Y = 96$~; somme 301.9~; coordonnées 31.8, 36.4, 31.8~: champ B. (b) Ti
16\,780, soit 167.8, 220, 87~; somme 474.8~; 35.3, 46.3, 18.3~: champ D. (c) Ti 4794, soit
47.9, 50, 60~; somme 157.9~; 30.3, 31.7, 38.0~: champ B. Le basalte (b) est un basalte intraplaque
(d'une île océanique ou d'un continent)~; (a) et (c), dans le champ B, peuvent être des MORB, des tholéiites
d'arc ou des basaltes calco-alcalins~: le diagramme ne peut pas trancher, et un autre (un diagramme
en araignée, l'anomalie en Nb) est nécessaire.
\end{solution}

\begin{solution}{exo:b2:magmatic-series:8}
(a) MORB $142/164 = 0.866$, \omterm{def:b2:elements-in-the-earth:morb-oib}{basaltes d'île océanique} $4/173 = 0.023$, basaltes d'arc $101/201 =
0.502$. (b) Avec des probabilités a priori égales, la règle de Bayes donne $0.866/(0.866 + 0.023 + 0.502) =
0.866/1.391 = 0.623$. (c) Il est très peu probable que le basalte soit un \omterm{def:b2:elements-in-the-earth:morb-oib}{basalte d'île océanique}
(probabilité $0.023/1.391 = 0.017$), mais on ne sait pas s'il s'agit d'un MORB ou d'un basalte d'arc~:
à peu près trois chances sur cinq contre deux.
\end{solution}

\begin{solution}{exo:b2:magmatic-series:9}
Moles pour \qty{100}{g}~: \ce{Na2O} $5.5/61.98 = 0.08874$, \ce{K2O} $4.5/94.20 =
0.04777$, \ce{Al2O3} $10.5/101.96 = 0.10298$. (a) $(0.08874 + 0.04777)/0.10298 =
0.13651/0.10298 = 1.326 > 1$~: la rhyolite est \emph{hyperalcaline}. (b) Le feldspath peut
prendre $0.10298$ des $0.13651$ moles d'oxyde alcalin~; l'excédent, $0.03353$, représente
$0.03353/0.13651 = 0.246$, soit un quart des alcalins, qui doivent entrer dans d'autres minéraux.
\end{solution}

\begin{solution}{exo:b2:magmatic-series:10}
(a) $s = (\lambda_F - \lambda_M)/(\lambda_F - \lambda_A)$~: pour $\lambda_F = 0.5$, $(0.5
- 0.8)/(0.5 - 0.4) = -3.0$~; pour $0.7$, $(-0.1)/(0.3) = -0.333$. (b) En résolvant
$s(\lambda_F - 0.4) = \lambda_F - 0.8$, on obtient $\lambda_F = (0.4s - 0.8)/(s - 1)$~:
Islande, $s = -1.10$~: $(-0.44 - 0.8)/(-2.10) = 0.590$~; Cascades, $s = -0.72$~:
$(-0.288 - 0.8)/(-1.72) = 0.633$. (c) Le \omterm{def:b2:magma-emplacement:magma}{magma} des Cascades, avec le $\lambda_F$ le plus grand,
perd son fer plus vite~: la magnétite l'en extrait. (d) $\ln(F/M)$ croît au rythme
$\lambda_M - \lambda_F$ et $\ln(A/F)$ au rythme $\lambda_F - \lambda_A$~; donc, pour chaque unité
gagnée par $\ln(A/F)$, $\ln(F/M)$ gagne $-s$~: 1.10 en Islande, 0.72 aux Cascades. $F/M$ monte
le long des deux tendances, mais plus vite dans la
tholéiitique, la tendance de Fenner.
\end{solution}

\begin{solution}{exo:b2:magmatic-series:11}
(a) BHVO-2, moles pour \qty{100}{g}~: \ce{SiO2} 0.82557, \ce{Al2O3} 0.13182, \ce{FeO}
0.15521, \ce{MgO} 0.18007, \ce{CaO} 0.20328, \ce{Na2O} 0.03580, \ce{K2O} 0.00545.
Orthose 0.00545 et albite 0.03580~; \ce{Al2O3} restant $0.13182 - 0.04125 = 0.09057$,
tout dans l'anorthite avec autant de \ce{CaO}~; \ce{CaO} restant $0.11271$ dans le diopside~;
$(\ce{Mg,Fe})\ce{O}$ restant $0.33528 - 0.11271 = 0.22257$ dans l'hypersthène. Silice
nécessaire~: $6\times0.00545 + 6\times0.03580 + 2\times0.09057 + 2\times0.11271 + 0.22257 =
0.87663$~; disponible 0.82557~; manque 0.05106. Chaque $(\ce{Mg,Fe})\ce{O}$ transformé de
l'hypersthène en olivine économise une demi-silice, au plus $0.22257/2 = 0.11129$~: en
transformer $0.10212$ suffit, et $0.12045$ reste dans l'hypersthène. La norme contient
olivine et hypersthène, ni quartz ni néphéline~: la roche est saturée en silice au sens
de la définition, une \emph{tholéiite à olivine} normative en hypersthène, de la \omterm{def:b2:magmatic-series:series}{série
tholéiitique}, comme l'édifie le stade bouclier d'un volcan hawaïen. (b) Moles~: 0.75732, 0.14712, 0.16008, 0.21092, 0.17832, 0.06131,
0.01592. Norme~: or 0.01592, ab 0.06131, an 0.06988, di 0.10843, hy $0.37100 - 0.10843 =
0.26257$~; silice nécessaire $0.09552 + 0.36786 + 0.13976 + 0.21686 + 0.26257 = 1.08257$~;
disponible 0.75732~; manque 0.32525. Transformer tout l'hypersthène en olivine économise
0.13129, il reste 0.19396~; chaque mole de \ce{Na2O} passée de l'albite à la néphéline
libère 4 \ce{SiO2}, donc $0.19396/4 = 0.04849$
des 0.06131 moles de \ce{Na2O} de l'albite deviennent néphéline. Néphéline dans la norme~: la
roche est sous-saturée en silice, un \emph{basalte alcalin}, de la \omterm{def:b2:magmatic-series:series}{série alcaline}.
\end{solution}

\begin{solution}{exo:b2:magmatic-series:12}
(a) Basaltes d'arc $116/141 = 0.823$, MORB $21/101 = 0.208$, \omterm{def:b2:elements-in-the-earth:morb-oib}{basaltes d'île océanique}
$3/137 = 0.022$. (b) $0.823/(0.823 + 0.208 + 0.022) = 0.823/1.053 = 0.782$.
(c) Numérateur $0.1\times0.823 = 0.0823$~; dénominateur $0.0823 + 0.6\times0.208 +
0.3\times0.022 = 0.0823 + 0.1248 + 0.0066 = 0.2137$~; probabilité $0.385$. (d) Le même test sur la même roche
donne 0.78 ou 0.39 selon la fréquence de chaque contexte là où la roche a été trouvée~: un \omterm{def:b2:magmatic-series:discrimination}{diagramme de
discrimination} donne une vraisemblance, et
le contexte géologique doit fournir le reste.
\end{solution}

\begin{solution}{pb:b2:magmatic-series:1}
\noindent\textbf{Partie I.}

\textbf{1.} Oxydes~: lave 1 99.17, lave 2 98.65, lave 3 98.20~; avec les volatils
99.97, 99.80 et 99.45, tous proches de 100. L'eau (0.70, 1.10, 1.20) est inférieure à
\qty{2}{\percent} et le dioxyde de carbone (0.10, 0.05, 0.05) inférieur à \qty{0.5}{\percent}~: les
trois sont acceptables.

\textbf{2.} Facteur $100/99.17 = 1.00837$~: \ce{SiO2} 50.82, \ce{Na2O} 2.92, \ce{K2O}
0.63, alcalins 3.55.

\textbf{3.} Lave 2~: $100/98.65 = 1.01368$, \ce{SiO2} 59.40, \ce{Na2O} 3.95, \ce{K2O}
1.17, alcalins 5.12. Lave 3~: $100/98.20 = 1.01833$, \ce{SiO2} 66.50, \ce{Na2O} 4.48,
\ce{K2O} 1.88, alcalins 6.36.

\textbf{4.} Lave 1 (50.82, 3.55)~: basalte. Lave 2 (59.40, 5.12~; la limite supérieure de O2 y est à
$5.9 + 1.1\times2.40/6 = 6.34$)~: andésite. Lave 3 (66.50, 6.36~; entre 63 et la droite de (69, 8)
à (77, 0), sous la droite de (63, 7) à (69, 8))~: dacite.

\textbf{5.} Analyses brutes~: (50.4, 3.52), (58.6, 5.05), (65.3, 6.25)~: les mêmes champs. Le
recalcul déplace chaque point d'au plus \qty{1.2}{\percent} de silice et ne change
rien ici, contrairement à l'\cref{exo:b2:magmatic-series:2}.

\textbf{6.} Le \omterm{def:b2:magmatic-series:qapf}{diagramme QAPF}, à partir du mode mesuré. Une roche plutonique de composition andésitique est riche
en plagioclase, avec peu de feldspath alcalin ou de quartz~: elle doit tomber dans le champ 10 ou 10*,
une diorite ou une diorite quartzifère (son mode décide).

\noindent\textbf{Partie II.}

\textbf{7.} Lave 1~: droites à 0.43 et 1.47, $\ce{K2O} = 0.63$~: moyen-K. Lave 2~: 0.81
et 2.28, 1.17~: moyen-K. Lave 3~: 1.13 et 2.96, 1.88~: moyen-K.

\textbf{8.} \ce{FeO_T} et \ce{MgO} exempts de volatils~: 8.67 et 7.97~; 6.18 et 3.85~; 4.07
et 1.63. Lave 1~: $3.55 + 8.67 + 7.97 = 20.19$~; $A = 17.6$, $F = 42.9$, $M = 39.5$.
Lave 2~: 15.15~; 33.8, 40.8, 25.4. Lave 3~: 12.06~; 52.7, 33.7, 13.5.

\textbf{9.} Lave 1, $A = 17.6$~: la droite entre (42.8, 17.1) et (46.9, 18.1) est à
$F = 42.8 + 4.1\times0.5/1.0 = 44.9$~; $42.9 < 44.9$, côté calco-alcalin, près de la
droite. Lave 2, $A = 33.8$~: $F = 49.6 + 2.9\times1.2/5 = 50.3$~; $40.8 < 50.3$,
calco-alcaline. Lave 3, $A = 52.7$~: $F = 40.8 - 2.9\times2.7/5 = 39.2$~; $33.7 < 39.2$,
calco-alcaline.

\textbf{10.} $\ce{FeO_T}/\ce{MgO}$~: $8.67/7.97 = 1.09$, $6.18/3.85 = 1.61$, $4.07/1.63 =
2.50$. Le rapport croît, comme dans tout basalte en différenciation, mais le fer lui-même
passe de 8.67 à 4.07, et $F$ diminue sur le triangle de 42.9 à 33.7~: pas d'enrichissement en fer. Une suite
tholéiitique monterait d'abord vers F, comme le fait le ferrobasalte de
\cref{exo:b2:magmatic-series:5}.

\textbf{11.} La \omterm{def:b2:magmatic-series:series}{série calco-alcaline}. La magnétite, qui cristallise dans un \omterm{def:b2:magma-emplacement:magma}{magma} oxydé,
a retiré le fer avec le magnésium~; les espèces oxydées proviennent de la plaque plongeante
qui se déshydrate dans une zone de subduction.

\noindent\textbf{Partie III.}

\textbf{12.} $\mathrm{Nb_N} = 2.0/0.658 = 3.04$, $\mathrm{La_N} = 7/0.648 = 10.80$,
$\mathrm{Th_N} = 1.1/0.0795 = 13.84$.

\textbf{13.} $3.04/10.80 = 0.281$, nettement en dessous des médianes des MORB (0.956) et des
\omterm{def:b2:elements-in-the-earth:morb-oib}{basaltes d'île océanique} (1.215), et même sous la médiane des basaltes d'arc (0.484)~: un creux
profond en niobium.

\textbf{14.} $13.84/3.04 = 4.55$~; le MORB médian du diagramme en araignée a $2.64/6.08 =
0.43$. Le thorium, comme les autres éléments enrichis dans les basaltes d'arc, a été apporté par la
plaque plongeante.

\textbf{15.} \ce{TiO2} exempt de volatils $0.95\times1.00837 = 0.958$~; $\mathrm{Ti} =
0.958\times5993 = \qty{5741}{ppm}$~; Ti/100 = 57.4, Zr 70, $3\mathrm Y = 54$, somme 181.4~:
31.6, 38.6, 29.8, dans le champ B, partagé par les MORB, les tholéiites d'arc et les basaltes calco-alcalins.
Il exclut un basalte intraplaque mais ne tranche pas entre une dorsale et un arc. Le
diagramme a été construit pour des basaltes~: dans une andésite et une dacite, le titane a été éliminé
par les oxydes et le zirconium concentré par la cristallisation, et leurs champs (la lave 2
tomberait dans C, la lave 3 hors de tous les champs) ne voudraient rien dire.

\textbf{16.} Une \omterm{def:b2:magmatic-series:series}{série calco-alcaline}, de type moyen-K, avec un creux profond en niobium et un
$\mathrm{Th_N}/\mathrm{Nb_N}$ élevé, sur un volcan qui a émis du basalte jusqu'à la dacite~: le volcan
s'est édifié dans un arc volcanique au-dessus d'une zone de subduction.

\noindent\textbf{Partie IV.}

\textbf{17.} $C_L = C_0/F$ pour un élément parfaitement incompatible~: $F = 0.63/1.17 =
0.538$.

\textbf{18.} Rubidium $10/20 = 0.500$~; thorium $1.1/2.1 = 0.524$. (Les \omterm{def:b2:elements-in-the-earth:trace}{éléments en
traces}, tels qu'analysés, changeraient de moins de \qty{1}{\percent} avec
le recalcul.)

\textbf{19.} $C_L/C_0 = F^{D - 1}$~: $125/70 = 1.786 = F^{-0.9}$, $F = 1.786^{-1/0.9} =
0.525$.

\textbf{20.} À partir de la potasse, $F = 1.17/1.88 = 0.622$ de la lave 2 à la lave 3. Pour le strontium~:
$380/480 = 0.792 = 0.622^{D - 1}$, donc $D - 1 = \ln 0.792/\ln 0.622 = -0.234/-0.474 =
0.49$ et $D \approx 1.49$~: le strontium est devenu compatible. Le plagioclase, qui fixe le strontium,
cristallisait en abondance.

\textbf{21.} À partir de la potasse, $F = 0.63/1.88 = 0.335$ (rubidium $10/32 = 0.31$, thorium
$1.1/3.4 = 0.32$)~: environ deux tiers du basalte, $1 - 0.335 = 0.665$, ont cristallisé
avant la dacite.

\textbf{22.} L'apport de potasse par assimilation élèverait \ce{K2O} dans la lave 2 au-dessus de
ce que la seule cristallisation donne, de sorte que $F = 0.63/1.17$ sortirait trop bas, et la masse de
cristaux trop élevée.

\textbf{23.} $(0.538 + 0.500 + 0.524 + 0.525)/4 = 2.087/4 = 0.522$~; cristaux $1 - 0.522
= 0.478$ du basalte, soit \textbf{\qty[text-series-to-math=true]{48}{kg} de
cristaux par \qty{100}{kg} de basalte} pour produire l'andésite.
\end{solution}
